\chapter{Literature Review}
\label{chap:litreview}

\section{The five models}

Goetz et al. treat threads and thread pools as the basic tools of concurrent Java \cite{ref1}. A fixed pool is both a scheduler and a limit: when every worker is blocked in I/O, further tasks queue and latency grows even if the dependency could accept more calls. The book does not treat a dependency’s own cap as a separate resource.

JEP 444 defines a virtual thread as a JVM thread that unmounts from its carrier during blocking I/O \cite{ref2}. Many blocked requests then occupy few OS threads, so one would expect the dependency cap or memory, not the pool, to be exhausted first; the JEP does not measure this. The Java 21 guide puts the limit in a \texttt{Semaphore} around the call to the capped service, and says that virtual threads should not be pooled for this purpose \cite{ref3}. \texttt{StructuredTaskScope} cancels unfinished subtasks when the scope ends, but it is still a preview API \cite{ref4}, so the Java baseline in this thesis is plain blocking code, with cancellation by thread interruption. These texts state the rule. They do not measure what happens at the cap, and they do not cover other models.

In Scala, a \texttt{Future} runs on the pool chosen by its \texttt{ExecutionContext}, does not limit how many tasks are submitted (the pool size limits only how many run at once), and a blocking call inside it holds a pool thread \cite{ref5}. ZIO schedules fibers on a worker pool, has a separate path for blocking work, and treats interruption as part of the fiber \cite{ref6}. In Kotlin, a suspend function frees its thread at a suspension point, dispatchers place coroutines on pools, and \texttt{coroutineScope} cancels the children together with the parent \cite{ref7}, a rule introduced by Elizarov in 2018 \cite{ref8}. A blocking call on a limited dispatcher still holds a thread. The documentation of these three describes the programming model. It does not separate the worker-pool size from an external cap, and it does not say whether a slot is released after cancellation.

\section{Saturation and where the limit goes}

Kleppmann describes how queueing delay and head-of-line blocking dominate the high percentiles of response time \cite{ref9}. Dean and Barroso show that for a request that waits on several components these percentiles decide what the user sees \cite{ref10}. Little’s law links the average number of requests in a system to their arrival rate and the time each spends there, $L = \lambda W$ \cite{ref11}. For a stub with $N$ slots and delay $D$, every call holds a slot for $D$, so $N$ slots serve at most $N/D$ calls per unit of time, and reach that rate only if callers keep all slots busy. The SRE book lists saturation among the signals to watch \cite{ref12}. Other sources say where a limit goes. The HikariCP note argues that once the database is the bottleneck, a larger connection pool adds contention, not throughput \cite{ref13}. Nygard lists bulkheads, timeouts, back pressure, and load shedding as patterns that bound work \cite{ref14}. Netflix computes the limit of in-flight requests at run time and publishes a Java library \cite{ref15}. All of them show that a limit must exist and why. None says which of the five models reaches the limit of the stub and which is stopped earlier by its own pool, and none checks whether a cancelled request returns its slot.

Two reports measure the placement on Java threads. Swift describes a database test in which a fixed pool of 40 platform threads limited concurrency by accident. Replacing it with virtual threads removed the limit and raised tail latency by an order of magnitude; a \texttt{Semaphore} with 20 permits brought the 99.999th percentile from over 10 s to 200 ms \cite{ref16}. Popov runs Spring Boot with virtual threads against a stub with a 200 ms delay and a cap of 1000 calls, and compares a limit at the server, in the application, and at the stub only. At 3x of capacity the server-level limit keeps 88% of peak goodput, the application-level limit about 70%, and with only the stub limit goodput collapses between 1.25x and 1.5x \cite{ref17}. Popov lists as open that a client timeout does not stop the work already running in the service. Both reports cover Java threads only. For virtual threads the placement is therefore partly answered; for \texttt{Future}, ZIO, and Kotlin on a blocking call it is not.

\section{Prior comparisons and the gap}

Comparisons of these models exist, but only in pairs and triples. Several sources are theses, blog posts, or mailing-list reports, so they are used as reported observations. For several academic items only the abstract was read.

\textbf{Java threads.} Lasic et al. compare platform and virtual threads in database-backed servers \cite{ref18}. Pandita finds equal throughput on CPU-bound work and higher throughput for virtual threads on I/O-bound work \cite{ref19}. Gustafsson and Nederlund Persson add reactive streams and limited and unlimited platform pools; virtual threads give lower latency and reactive streams lower memory use \cite{ref20}. Hu is the closest. In SPECjbb2015 on JDK 25 virtual threads have lower peak throughput, and bounding their number raises the throughput metrics by about 16--20% \cite{ref21}. This supports the claim that an explicit bound matters, but the study covers only the two thread types and no capped dependency.

\textbf{Virtual threads and Kotlin coroutines.} Beronić et al. compare platform threads, virtual threads, and coroutines on structured-concurrency constructs \cite{ref22}. Grabowiec et al. find coroutines on the IO dispatcher about 16% faster than a thread pool on a REST workload, with virtual threads behind at higher iteration counts \cite{ref23}. Mohammed compares virtual threads and coroutines on four JVMs: virtual threads scale better on I/O, coroutines have lower overhead on CPU work \cite{ref24}.

\textbf{Effect systems.} Spiewak shows that the ZIO and Cats Effect microbenchmarks each stress a narrow artificial workload \cite{ref25}. VirtusLab compares \texttt{Future}, Cats Effect, ZIO, Kyo, Gears, and Ox on a crawler task with a parallelism limit and cancellation on failure, as a code comparison without measurement \cite{ref26}. A blog post tabulates platform threads, virtual threads, ZIO fibres, and Kotlin coroutines, including interruption, with approximate figures \cite{ref27}. It is the nearest published property table, but it has no \texttt{Future}, no dependency cap, and no measured cancellation. Queries in OpenAlex, Crossref, and arXiv returned no peer-reviewed comparison of ZIO or Cats Effect with other concurrency models.

\textbf{Cancellation.} Ward’s 2024 talk notes compare racing, loser cancellation, and resource closing in ZIO, Kyo, Ox, Java with Loom, and Rust, and record that loser cancellation was not validated in one scenario \cite{ref28}. AsyncLeakBench collects 902 fixed asynchronous resource leaks in Java projects, with categories triggered by cancellation and timeouts, and reports that detectors handle those triggers poorly \cite{ref29}. This shows that a resource can stay held after cancellation, but not for a slot in a concurrency limit.

\textbf{Benchmark suites.} Renaissance, the peer-reviewed suite of concurrent JVM workloads, was built to evaluate compilers and VMs \cite{ref30}. Its workloads include Scala \texttt{Future}, Akka, Rx, Finagle, and Ktor, and it has no capped dependency, no request-level tail latency, and no cancellation scenario \cite{ref31}. It also advises that a few repetitions give only a quick look and a thorough evaluation needs many more \cite{ref31}; the benchmark chapter follows this.

\textbf{What is not covered.} In the sources found, no single study combines:

\begin{enumerate}
\item all five models of this thesis, with \texttt{Future} and ZIO next to Java threads and Kotlin;
\item the same request in every model, with a call to an external dependency that has its own cap;
\item the dependency cap and the pool size varied separately, to identify the first exhausted resource;
\item measured cancellation, including whether a slot in the cap is released when a fan-out is cancelled.
\end{enumerate}


\vspace{0.5em}\noindent\rule{\textwidth}{0.5pt}\vspace{0.5em}

